\addcontentsline{toc}{subsubsection}{Solution 49. The Prime-Pairing Machine}

\begin{dlsummary}[title={\bfseries Solution 49. The Prime-Pairing Machine},breakable]

\begin{proof}
\textbf{The solutions:} Partition the set of all primes into pairs. Within each pair, designate one prime \(p\) and the other \(q\), and set
\[
\boxed{f(p)=q,\qquad f(q)=p^2.}
\]
Extend these assignments multiplicatively to all positive rationals. Every solution arises this way.

\paragraph{1. Remove the inequality.}

Positivity and multiplicativity give
\[
f(1)=1,\qquad f(1/x)=\frac1{f(x)}.
\]
Applying the composition inequality to \(1/x\),
\[
\frac1{f(f(x))}\ge\frac1{x^2}.
\]
Together with the original inequality, this proves
\[
\boxed{f(f(x))=x^2\qquad(x>0).} \tag{1}
\]

\paragraph{2. Show that primes map to prime powers.}

Let \(p\) be prime. Since \(f(p)\) is an integer and
\[
f(f(p))=p^2>1,
\]
we have \(f(p)>1\).

Choose a prime divisor \(q\) of \(f(p)\). Factoring \(f(p)\) into primes and applying \(f\), all factors are positive integers and their product is \(p^2\). Also, \(f(q)>1\), by the same reasoning as above. Thus
\[
f(q)=p^b
\]
for some positive integer \(b\).

Apply \(f\) once more:
\[
q^2=f(f(q))=f(p)^b.
\]
Unique prime factorisation forces
\[
\boxed{f(p)=q^a}
\]
for a positive integer \(a\). In particular, \(f(p)\) has only one prime divisor.

\paragraph{3. Identify the allowed pairs.}

Now
\[
p^2=f(f(p))=f(q)^a=p^{ab},
\]
so
\[
ab=2.
\]
Therefore
\[
(a,b)=(1,2)\quad\text{or}\quad(2,1).
\]

Moreover, \(q\ne p\): otherwise \(f(p)=p^a\) would give \(a^2=2\), impossible for an integer \(a\).

Thus every prime belongs to a two-prime pair, with assignments
\[
f(p)=q,\qquad f(q)=p^2
\]
after choosing the orientation. These pairs are disjoint because the prime associated with any given prime is uniquely determined.

\paragraph{4. Construct and verify every solution.}

Choose any partition of the primes into oriented pairs, make the assignments above, and define
\[
f\!\left(\prod_r r^{e_r}\right)=\prod_r f(r)^{e_r},
\]
where the exponents \(e_r\) are integers and only finitely many are nonzero.

This is well-defined and multiplicative. It sends positive integers to positive integers, and on every prime,
\[
f(f(r))=r^2.
\]
Multiplicativity therefore gives \(f(f(x))=x^2\) for every positive rational \(x\), verifying all conditions.

\paragraph{Example.}

Pair consecutive primes and orient them as
\[
(2,3),\ (5,7),\ (11,13),\ldots.
\]
Then
\[
f(2)=3,\quad f(3)=4,\quad f(5)=7,\quad f(7)=25,
\]
so, for instance,
\[
f(6)=12,\qquad f(12)=36.
\]
This construction also passed direct composition checks on the integers from 1 through 100.[1]
\end{proof}
\end{dlsummary}

\begin{takeaways}
\begin{itemize}
\item \textbf{Reciprocals can force exactness:} A global multiplicative inequality may secretly be an equation.
\item \textbf{Integrality controls factorisation:} Integer-valued images prevent unwanted prime factors from cancelling.
\item \textbf{Follow a prime through two iterations:} The return to \(p^2\) restricts both the prime support and the exponents.
\item \textbf{Classification need not mean uniqueness:} A complete answer can be a structural recipe for infinitely many functions.
\end{itemize}
\end{takeaways}
